\scanpage{148}
c) Soit $Q=F^\perp/F$, alors $q$ induit par passage au quotient une forme \deleted{\unsure{Q}} $q_Q$ sur $Q$, et \deleted{la forme bil. associée est} \inserted{celle-ci est régulière}

d) Soit
\[
\theta:G\longrightarrow E
\]
un scindage de la suite exacte
\[
0\mathbin{-}F^\perp\longrightarrow E\longrightarrow G\longrightarrow0\tag*{($*$)}
\]
\deleted{\unsure{et}} et considérons la forme quadratique
\[
q_\theta:F\times G\longrightarrow\underline{\mathcal O}_S
\]
\unsure{définie|défini} par
\[
q_\theta(x,y)=q(x+\deleted{\unsure{\varphi}}\theta(y))=q(\theta(y))+\deleted{\unsure{\varphi}}\langle y,x\rangle
\]
La forme bilinéaire associée est
\[
\begin{aligned}
\varphi_\theta(x,y\,;x',y')&=\varphi(x+\theta y,\theta y')+\varphi(x',\theta y)\\
&\inserted{=\varphi(\theta y,\theta y')+\langle y',x\rangle+\langle y,x'\rangle}
\end{aligned}
\]
et elle est non \uline{dégénérée}, donc $q_\theta$ est une forme quadratique régulière (de rg pair). Si $Q_\theta$ est l'orthogonal de $F+\theta(G)$, on a \deleted{donc}
\[
E\simeq F\oplus Q_\theta+\theta(G)
\]
\inserted{et \deleted{\unsure{$Q_\theta$ forme|la forme}} \inserted{la forme} $q|Q_\theta$ est régulière. $[$\uline{NB} $Q_\theta\simeq Q=F^\perp/F$, et la forme $q_Q$ déduite de $q_\theta$ ne dépend pas du choix de $\theta$, étant obtenue par passage au quotient\ldots}%
% slanted

d) Supposons $S$ affine, \deleted{\unsure{s}} ce qui assure l'existence d'un $\theta\ldots$ \deleted{\unsure{s}}. Alors il existe (pour $\theta$ fixé)
\[
u:G\longrightarrow F
\]
tel que, posant
\[
\theta_u^{\deleted{\unsure{s}}}=\theta+u
\]
on ait
\[
q\circ\theta_u=0
\]
\margin{\uline{NB} condition qui équivaut à $(q\circ\theta)(y)+\langle y,u(y)\rangle=0$ i.e. $B_u(y,y)=-(q\circ\theta)(y)$ où $B_u$ est la forme bil. associée à $u$}
i.e. $\theta_u(G)$ soit tot. isotrope dans $E$ (scindage isotrope). L'ensemble des \deleted{\unsure{sections (?)|$u\in\mathrm{Hom}($}}
\[
u\in\mathrm{Hom}(G,\underset{\substack{\shortparallel\\\check G}}{F})\simeq\mathrm{Bil}(G,G)
\]
tels que \deleted{$q\circ$} $\theta_u$ soit isotrope, est un torseur sous le sous-groupe $\mathrm{Alt}(G,G)$ de $\mathrm{Bil}(G,G)$.

Ainsi (indépendamment du choix d'un $\theta$)
% archive sheet number: 147
